\documentclass[11pt,a4paper]{article}
\usepackage[hmargin=0.65in,vmargin=0.55in]{geometry}
\usepackage[T1]{fontenc}
\usepackage{lmodern}
\usepackage{microtype}
\usepackage{hyperref}
\usepackage{titlesec}
\hypersetup{colorlinks=true,urlcolor=blue}
\setlength{\parindent}{0pt}
\setlength{\parskip}{1pt}
\titleformat{\section}{\normalsize\bfseries}{}{0pt}{}[\titlerule]
\titlespacing*{\section}{0pt}{5pt}{3pt}

\begin{document}
\begin{center}
{\large\bfseries Research proposal and eligibility statement}\\[3pt]
Phurailatpam Hemantakumar\\
Juan de la Cierva 2026 --- IGFAE, Universidade de Santiago de Compostela
\end{center}

\section*{Research proposal}
At IGFAE, I would develop gravitational-wave (GW) lensing statistics and waveform models to infer compact-binary populations, environments, and cosmology. My experience combines population simulation, analytical waveform development, scientific software, and Bayesian data analysis, matching IGFAE's interests in gravitational waves, star formation, and cosmology.

\textbf{Strong-lensing statistics and LVK commitments.}
Assessing lensing candidates requires reliable rates and predictions of observable populations. I lead development of \texttt{ler} for population simulations and developed \texttt{gwsnr} for rapid SNR and detector-selection calculations. I have implemented propagation of source and lens model uncertainties and finite Monte Carlo sampling errors into lensing statistics, including detectable time-delay distributions. As project leader, I will continue my ISR-weighted LVK/GRASP commitments for \texttt{ler} towards GWTC-6 and GWTC-7: improved rate forecasts, magnification and time-delay priors, and four-image and near-caustic forecasts including missing images. With Chunglee Kim and Jiyun Park, I will extend our ongoing star-formation-informed merger-rate work; with Anupreeta More, I will extend predictions to galaxy-cluster lenses; injection-based sensitivity calibration in \texttt{gwsnr} will improve detector selection. I will test how these predictions affect candidate odds and maintain the software for LVK analyses.

\textbf{Time-delay cosmography.}
The detectable lensed-event fraction and time-delay distribution carry cosmological information, but also depend on source and lens populations and selection. Building on \href{https://arxiv.org/abs/2512.15168}{Maity et al.}, I propose to use my existing uncertainty-propagation framework to test the robustness of this method. I will compare alternative source populations and lens mass models, propagate their uncertainties into cosmological inference, and test bias and credible-interval coverage with simulated catalogues. This will identify constraints robust to astrophysical assumptions and the population information needed to improve them.

\textbf{Eccentric waveforms and inference.}
Residual eccentricity probes binary formation, but waveform errors can bias its interpretation. I co-developed \texttt{TaylorF2Ecck}, implemented it in a \texttt{LALSuite} fork, and analysed GW170817 and GW190425. Our studies identified mass-ratio biases from insufficient phasing accuracy. Building on ongoing spin extensions by Shri Krishna Anand, guided by Achamveedu Gopakumar, I will implement and validate spin effects, higher-order eccentricity terms, and amplitude corrections, and pursue eccentric phasing through at least 3.5PN order. Injection studies will establish when these improvements support reliable parameter recovery. I plan to add this waveform development and analysis to my GRASP pledges and pursue integration into the official \texttt{LALSuite} repository.

\textbf{Line-of-sight acceleration (LOSA).}
Our GW190814 study found a strong LOSA--eccentricity degeneracy without significant evidence for either effect. I propose a follow-up using improved eccentric waveforms and \href{https://doi.org/10.1103/swkb-nkgg}{Zhao et al.}'s frequency-domain spectral-differentiation method to model acceleration-induced time shifts without the stationary-phase approximation. Joint inference and injection tests across more events will assess whether the effects can be distinguished, enabling more reliable tests of binary environments and possible tertiary companions.

\textbf{Multimessenger lensing.}
With Otto Akseli Hannuksela, Albert Kong and other collaborators, I will complete and validate our ongoing Bayesian lensing tests for multiple GRBs, then incorporate light-curve information and explore joint GW--GRB associations.

These projects build on my completed and ongoing research and existing collaborations.

\section*{Eligibility for Juan de la Cierva 2026}
I defended and received approval of my doctoral thesis in Physics at The Chinese University of Hong Kong on 30 August 2026, under the supervision of Otto Akseli Hannuksela. My formal PhD degree certificate is pending issuance. Article 6 of the 2026 call defines the date of obtaining the doctorate as the date of the thesis defence and approval; my defence date falls within the required period of 1 January 2024--31 December 2026. My doctoral research was carried out principally at The Chinese University of Hong Kong, which is distinct from the proposed host institution, IGFAE at the Universidade de Santiago de Compostela. I have not applied to the 2026 Ram\'{o}n y Cajal call and have not previously been awarded a fellowship under Juan de la Cierva, Formaci\'{o}n Posdoctoral, Juan de la Cierva-formaci\'{o}n, or Juan de la Cierva-incorporaci\'{o}n. I did not participate in the 2025 Juan de la Cierva call.

% Documentation: For a doctorate from a non-Spanish university, the official
% 2026 call requires a copy of the doctoral degree or academic certification
% explicitly stating the date of award (the defence and approval date).
% Obtain CUHK certification of both the defence and its successful approval.
% Article 6(1)(e) does not apply: no participation in the 2025 call.
% Personal eligibility declarations confirmed by the applicant on 3 November 2026.

\end{document}
